本研究旨在解决激光超声波无损检测中的一个关键空白：多域融合神经网络是否能够在减少平均次数的条件下替代传统信号平均，同时保留对缺陷检测至关重要的声学特征。
维纳滤波和小波阈值等传统去噪方法能够抑制噪声，但频繁破坏声学有意义的特征——包括到达时间、幅度和波形形态\cite{chen2019wavelet}。与此同时，现有深度学习方法要么需要大量信号平均进行训练\cite{liu2022deep}，要么在单一域中操作，无法捕捉互补的时频信息\cite{zhang2020deep}。这一限制对于激光超声波无损检测尤为突出，因为重复平均会导致表面损伤并降低检测效率\cite{yang2023unsupervised}。本研究提出的核心问题是：\emph{在减少平均次数的条件下训练的多域融合去噪网络，能否相比传统信号平均实现相当或更优的缺陷检测能力，同时实现从未损伤样本到含缺陷工件的泛化？}

为回答此问题，本文做出以下创新贡献：

\begin{enumerate}
\item \textbf{多域融合架构。} 一种物理约束的编码器-解码器网络，联合编码时域波形和时频表示，充分利用两个域的互补信息以提高信号增强保真度。

\item \textbf{物理约束损失函数。} 一种明确保留声学有意义特征的损失 formulation，包括到达时间、峰值幅度和波形形态，确保去噪后的信号保留用于缺陷表征的诊断信息。

\item \textbf{减少平均训练策略。} 一种配对训练策略，其中网络输入为有限平均信号（$N$ 次平均），目标输出为高平均参考信号（$\gg N$），使模型能够在实际减少平均条件下学习去噪。

\item \textbf{跨域泛化评估。} 系统性验证表明，仅在完整铝试样上训练的模型在含缺陷工件上保持检测性能，解决了先前深度学习无损检测研究中识别的从无缺陷到有缺陷的泛化差距 \cite{liu2022deep}。
\end{enumerate}

这些贡献共同建立一个实用激光超声波检测框架，减少了对耗时的信号平均的依赖，同时保留——而非牺牲——缺陷检测能力。